\documentclass[11pt]{article}
\usepackage[T1]{fontenc}
\usepackage{lmodern}
\usepackage[margin=1in]{geometry}
\usepackage{amsmath,amssymb,amsthm,booktabs,array,microtype}
\usepackage{tikz,pgfplots}
\pgfplotsset{compat=1.18}
\usepackage[colorlinks=true,linkcolor=blue!45!black,citecolor=blue!45!black,urlcolor=blue!45!black]{hyperref}
\newtheorem{proposition}{Proposition}
\newcommand{\R}{\mathbb{R}}
\newcommand{\norm}[1]{\left\lVert #1\right\rVert_2}
\newcommand{\code}[1]{\texttt{#1}}
\title{Cotangent: Preserving Reduction Arithmetic\\in Batched Gradient Handling}
\author{Atila Vahedian}
\date{October 2026}
\begin{document}
\maketitle
\begin{abstract}
Gradient handling can occupy a substantial fraction of small-model training on
an accelerator even when the optimizer is already fused. Batching this work is
straightforward in real arithmetic but can change floating-point reductions.
We study a gradient layout that preserves the native two-level clipping norm
while batching equal-length tensors on PyTorch's Metal backend. The method
retains every model derivative and applies native fused AdamW to consistently
packed parameter and gradient coordinates. An independently frozen experiment
with 100 paired seeds reduces time to a fixed validation target by 15.03\%
(99.5\% confidence interval, 13.50--16.56\%) relative to eager native fused
AdamW, while meeting a predeclared 0.01-bit-per-byte quality tolerance.
@@EXTENSION_ABSTRACT@@
Numerical checks cover complete model updates, optimizer states, reduction
boundaries and composition with the compiler. The results are specific to one
Apple M5 Pro, the pinned PyTorch backend and byte language models; the study
does not establish performance on other accelerators or model scales.
\end{abstract}

\section{Introduction}
The work after backward is easy to underestimate. A training step must combine
gradients to determine a clipping coefficient, scale them, update moment
estimates, and change parameters. A fused optimizer reduces part of this cost,
but the surrounding norm calculations and tensor dispatch can remain visible
when a model consists of many modestly sized parameter tensors.

Cotangent investigates that remaining cost. The starting point is a conventional
byte transformer with native attention, linear and embedding derivatives. Its
clipping operation evaluates 37 individual parameter norms before combining
them. The proposed layout batches those norms into four equal-length groups
and exposes one contiguous coordinate vector to native fused AdamW.

The difficulty is numerical. A flat global norm has the same mathematical value
as a norm of parameter norms, yet the two computations can use different
floating-point additions. Even preserving a separate row norm is insufficient
if tensor shape selects a different kernel. We therefore preserve the native
two-level norm structure and choose a layout that retains the pinned MPS generic
reduction path. Correctness is checked before timing is interpreted.

Two experiments serve different purposes. The first is a powered replication
with a fixed time-to-quality target and quality tolerance. The second extends
the evidence with a factorial ablation, additional model layouts, and stronger
compiler baselines. In the extension, both arms use the same selected model
execution path. This isolates the additional effect of gradient handling from
the benefit of compilation itself. Code, protocols, observations and analysis
are available at \url{https://github.com/atilavahedian/cotangent}.

\section{Method}
\subsection{Preserve the clipping computation}
Let $g_\ell\in\R^{K_\ell}$ be the completed gradient of parameter tensor $\ell$.
For a clipping threshold $\tau$, native clipping computes
\begin{equation}
 n_\ell=\norm{g_\ell},\qquad
 n=\left(\sum_{\ell=1}^{L}n_\ell^2\right)^{1/2},\qquad
 c=\min\left(1,\frac{\tau}{n+\epsilon_c}\right),\qquad h_\ell=cg_\ell,
 \label{eq:clip}
\end{equation}
where $\tau=1$ and $\epsilon_c=10^{-6}$ in these experiments. We retain a
separate norm for every parameter and retain the original order in the final
reduction. No gradient coordinate is sampled, dropped or approximated.

For a bucket containing $B$ gradients of length $K$, stack their flattened
coordinates and view the resulting tensor as $B\times(K/2)\times2$. Reducing
the two trailing axes produces $B$ independent parameter norms. For even
$K\ge4$, both reduced axes are nontrivial. Short or odd lengths use separate
native norms. Scalar outputs are restored to parameter order before the final
norm in Equation~\ref{eq:clip}.

This choice follows a specific dispatcher condition \cite{mps}. At the pinned PyTorch
revision, the MPS inner norm fast path requires multiple output elements and
exactly one nontrivial reduced dimension on the innermost axis. A conventional
$B\times K$ row reduction meets those conditions. The two-axis layout avoids
that path and invokes the generic norm kernel used by a native scalar parameter
norm. Each parameter keeps $K$ reduced elements and the same threads per group.
The source-level argument concerns this particular implementation; it is
supported by measured numerical checks, rather than asserted for arbitrary
backends or versions.

\subsection{Pack coordinates consistently}
The model retains independent autograd leaves. Before its first forward pass,
their storage is attached to disjoint views of a master parameter vector.
Backward still materializes every native parameter gradient. The completed
gradients are concatenated, multiplied by $c$, and passed to native fused
AdamW. The tensor copies, concatenation and temporary allocations are part of
the measured training time.

The supported setting has one hyperparameter group, homogeneous contiguous
parameters, dense gradients for all active coordinates, and independent,
nonoverlapping parameter storage. Parameters must not be moved or reassigned
after construction. Partial updates, mixed groups and sparse gradients require
additional treatment and are outside the implemented experiment.

\begin{proposition}[Coordinate equivalence in real arithmetic]
Let $P$ be the fixed reindexing that concatenates all active coordinates.
With identical initialization, hyperparameters, scalar step counts and complete
gradients, clipping followed by AdamW commutes with $P$ in real arithmetic.
\end{proposition}
\begin{proof}
The reindexing is an isometry, so $\norm{Pg}=\norm{g}$ and the scalar clipping
coefficient is unchanged. Let $h=cg$. AdamW's moments satisfy
\begin{align}
 m_t&=\beta_1m_{t-1}+(1-\beta_1)h_t,\\
 v_t&=\beta_2v_{t-1}+(1-\beta_2)h_t^2.
\end{align}
Reindexing commutes with coordinatewise addition, multiplication, squaring,
square roots and division. Bias correction uses the same scalar $t$.
Consequently the decoupled weight decay and adaptive update also commute with
$P$. Induction from the same initial parameters and zero moments gives
coordinate-equivalent iterates.
\end{proof}

This algebra does not establish bitwise equality of parallel reductions.
Floating-point addition is not associative. The implementation therefore keeps
the native reduction structure and tests the remaining finite-precision
behavior on the actual device. Those checks also distinguish optimizer
equivalence from run-to-run variation in independently executed native
embedding accumulation.

\section{Experimental design}
\subsection{Environment, models and data}
All GPU measurements use one Apple M5 Pro with 24 GiB RAM, macOS 27.0,
Python 3.12.14 and PyTorch 2.14.1 MPS. The Torch source revision is
\code{5c4886908584029761b579af026dcfb627c84070}.
Weights and completed gradients are FP32; forward computation uses BF16
autocast. AdamW uses a peak learning rate of $10^{-3}$, $\beta_1=0.9$,
$\beta_2=0.95$, $\epsilon=10^{-8}$, weight decay 0.1 and global clipping at 1.
The learning rate warms up and then follows a cosine schedule to one tenth of
the peak. Batch size is eight and sequence length is 256.

The corpus is a pinned byte serialization of WikiText-2 \cite{wikitext}.
The first 95\% of training bytes is used for optimization; the remaining 5\%
is a diagnostic split. Full official validation and test evaluations cover
1,148,006 and 1,292,012 predicted bytes respectively. These corpus splits were
exposed during development and are explicitly reused as quality regression
data. Fresh seed pairs do not make the corpus an unseen generalization test.
All models train from scratch and use an untied output head.

\begin{table}[t]
\centering
\small
\begin{tabular}{lrrrr}
\toprule
Model & Parameters & Width & Layers & Tensors / buckets\\
\midrule
@@ARCH_TABLE@@
\bottomrule
\end{tabular}
\caption{Model configurations in the extension. Bucket counts refer to
equal-length completed gradient tensors, including norm parameters.}
\label{tab:models}
\end{table}

The two original transformers use pre-norm attention and GELU feed-forward
layers. The additional decoder uses six query heads, two key/value heads,
RMSNorm and a SwiGLU-style gated feed-forward layer \cite{gqa,rms,glu}.
Its key/value heads are explicitly repeated for the native attention call.
The convolutional model uses causal, depthwise, kernel-three convolutions with
dilations $1,2,4,8,16,32$, residual pointwise mixing and GELU feed-forward layers,
following the causal dilated design of WaveNet \cite{wavenet}.
Architecture-level language modeling quality is not the comparison objective:
each optimizer treatment is compared within the same architecture.

\subsection{Independent time-to-quality replication}
The confirmatory experiment (V5) contains 100 paired seeds and 200 runs of the
3.35M transformer, each trained for 2,000 steps. Every pair matches
initialization, data schedule, architecture, precision and optimizer settings.
Exactly 50 blocks run native first and 50 run Cotangent first; block order is
shuffled before training. The protocol is public at tag \code{protocol-v5}.

The quality target is 3.2 bits per byte (BPB) on a scheduled validation probe
covering 65,536 targets. Probes occur every 100 steps. The target time is the
first observed crossing, with no interpolation. Its clock includes setup,
priming, transfer, forward/backward, gradient handling, optimizer updates,
equal host/resource checks and preceding probes. Python process launch is
excluded equally. Every training step synchronizes MPS.

The fixed decision requires at least 10\% mean paired target-time reduction,
a positive timing lower bound, and an upper full-test degradation bound no
greater than 0.01 BPB. Intervals are two-sided 99.5\% paired Student-$t$
intervals. The sample size was fixed before V5 outcomes using all 25 quality
differences from the preceding V4 study, whose standard deviation was
0.0209261 BPB. A normal planning approximation at zero degradation and 90\%
power gives 73.2 pairs; the protocol rounds up to 100. V4's failed combined
decision remains unchanged. V5 changes neither the method nor its thresholds.

\subsection{Ablations and stronger execution paths}
The extension (V6) is a separately frozen descriptive study with 128 runs.
Eight seed blocks cross scalar/bucketed norms with native/packed fused AdamW,
using 800 steps per cell on the small transformer. A rotating Latin order puts
each of the four cells in each execution position twice. This separates norm
batching, parameter packing and their interaction.

The breadth study uses twelve paired seeds per configuration, 1,200 steps per
run, and balanced six-AB/six-BA execution orders. It screens native fused,
automatic and loop AdamW, direct native multi-tensor norm/multiply operations,
TorchInductor, and AOT eager capture using two training-only diagnostic seeds.
Selection uses mean training time after the first 20 pilot steps.
TorchInductor is the fastest eligible baseline in all four configurations.
Compilation uses \code{fullgraph=True} and static shapes. Both official arms
therefore use the same compiled forward/backward; Cotangent changes the
subsequent norm/layout/update handling.

Compilation and setup are reported separately and are included in total
completion time. Process state is fresh for every run; generated compiler
artifacts use a shared cache already exercised by pilots. These are warm-cache
measurements, not a characterization of first-ever compilation latency.
Gradient handling and native fused AdamW remain outside the compiled model
call. The study does not exhaust every compiler/optimizer configuration.

All V6 blocks are shuffled before outcomes and the sample sizes are fixed.
V6 reports per-configuration descriptive 99\% paired Student-$t$ intervals;
it does not introduce a new powered quality gate or simultaneous coverage
across every comparison. Official validation/test data are evaluated once after
each run and are not used to select the execution paths.

For paired times $T_i^{\mathrm{base}}$ and $T_i^{\mathrm{cot}}$, the reported
reduction is the mean of $s_i=1-T_i^{\mathrm{cot}}/T_i^{\mathrm{base}}$.
It is not the ratio of aggregate mean times. For paired differences $z_i$, the
interval is $\bar z\pm t_{n-1,1-\alpha/2}s_z/\sqrt n$. Full-test quality
differences use candidate minus baseline BPB; positive values indicate worse
quality. All observations, including adverse pairs, are retained.

Sampled process RSS is capped at 5 GiB, MPS driver allocation at 8 GiB, and
swap growth at 512 MiB. Checks occur every ten steps. Training runs are serial;
no task-owned GPU jobs execute concurrently.

\section{Results}
\subsection{Numerical behavior}
The original check covers 32 full-model updates: two model sizes, FP32/BF16
forward precision and eight consecutive steps. Losses, gradients, clipping
norms, parameters and both moment buffers match bitwise. Seven additional
full-training derivative-controlled pairs finish with identical whole-model
hashes and full evaluation scores. Both arms use matched deterministic
embedding derivatives in these controls to isolate optimizer arithmetic.
Primary performance comparisons retain native embedding backward.

The extension adds 128 full-model checks across all four configurations,
two forward precisions, four gradient-handling treatments and four consecutive
updates. All checked losses, gradients, norms, parameters and moments match
the native clipping/fused AdamW reference. A further 16 BF16 checks compare
the two compiled paths with matched deterministic derivatives; losses,
gradients, norms and resulting parameter hashes match. These compiled checks
compare identical compiler paths, not compiled arithmetic with eager arithmetic.

A norm sweep checks 32 lengths, bucket sizes 1, 3 and 8, and FP32, FP16 and
BF16 gradient tensors: all 288 total norms match bitwise on this backend.
It includes short/odd fallback lengths and values near reduction-size boundaries.
Missing gradients, mixed parameter types and overlapping parameter objects are
rejected. These observations do not establish equivalence for every tensor
stride, nonfinite input, compiler transformation or future backend version.

\subsection{Confirmatory replication}
All 200 V5 runs complete and both arms reach the target in every pair.
Cotangent reduces target time by 15.03\% (99.5\% CI, 13.50--16.56\%) and
fixed-budget training time by 16.97\% (16.26--17.69\%). Average target times
are 16.65 seconds for native and 14.13 seconds for Cotangent; average training
times are 34.84 and 28.92 seconds. Ninety-nine pairs reach the target faster
with Cotangent; the one slower pair remains in the analysis.

Mean full-test degradation is 0.00138 BPB (99.5\% CI, $-0.00451$ to
$0.00728$). Its upper bound is below the predeclared 0.01 tolerance. This
supports the bounded noninferiority decision, rather than improved model
quality. Fifty of the 100 pairs have worse candidate test BPB.

\subsection{Factorial ablation}
@@FACTORIAL_TEXT@@
\begin{table}[t]
\centering
\begin{tabular}{lrr}
\toprule
Contrast & Training-time reduction (\%) & 99\% CI\\
\midrule
@@FACTORIAL_TABLE@@
\bottomrule
\end{tabular}
\caption{Eight paired factorial blocks, 800 steps per cell. Single-factor
contrasts compare with scalar norms/native layout. Main effects average the
factor's two conditional effects and normalize by the block's reference time.
The signed interaction is $(T_{00}-T_{10}-T_{01}+T_{11})/T_{00}$.}
\label{tab:factorial}
\end{table}

\subsection{Composition with compiled training}
@@BREADTH_TEXT@@
\begin{table}[t]
\centering
\small
\begin{tabular}{lrrr}
\toprule
Model & Base / Cot. time (s) & Reduction (\%) & 99\% CI\\
\midrule
@@BREADTH_TABLE@@
\bottomrule
\end{tabular}
\caption{Fixed-budget training, twelve pairs per model, 1,200 steps per arm.
Every baseline uses compiled forward/backward plus native clipping and fused
AdamW. Every candidate uses the same compiled model plus Cotangent handling.
Intervals are descriptive per configuration.}
\label{tab:breadth}
\end{table}

\begin{figure}[t]
\centering
\begin{tikzpicture}
\begin{axis}[width=.9\linewidth,height=5.2cm,xmin=0,xmax=@@PLOT_MAX@@,
 xlabel={Fixed-budget training-time reduction (\%)},
 ytick={1,2,3,4},yticklabels={Small transformer,Medium transformer,Gated GQA,Causal convolution},
 ymin=.5,ymax=4.5,grid=major,major grid style={gray!15},
 tick label style={font=\small},label style={font=\small},
 error bars/x dir=both,error bars/x explicit]
@@BREADTH_PLOT@@
\end{axis}
\end{tikzpicture}
\caption{Additional benefit of gradient handling on top of the same compiled
model. Points are means of paired reductions; horizontal bars are descriptive
99\% confidence intervals.}
\label{fig:breadth}
\end{figure}

\begin{table}[t]
\centering
\small
\begin{tabular}{lrrr}
\toprule
Model & Base / Cot. test BPB & Difference & 99\% CI\\
\midrule
@@QUALITY_TABLE@@
\bottomrule
\end{tabular}
\caption{Full-test quality in the extension. Difference is candidate minus
baseline BPB; positive is worse. These descriptive intervals do not replace
the independent V5 quality decision.}
\label{tab:quality}
\end{table}

@@QUALITY_TEXT@@
@@RESOURCE_TEXT@@
Every one of the 128 extension snapshots is restored locally. Weight-file
hashes, reconstructed parameter hashes and finite CPU inference are verified.
This audit does not repeat the official test evaluation. All weights remain
local; public records contain the hashes needed to identify the measured runs.

\section{Related work}
AdamW's coordinatewise moment and decoupled-decay updates provide the optimizer
semantics used here \cite{adamw}. Parameter flattening is also a standard
training-systems mechanism, including PyTorch FSDP \cite{fsdp}. DeepSpeed
FusedAdam and NVIDIA Apex's multi-tensor norm implementation batch work across
parameter tensors \cite{deepspeed,apex}. The present study evaluates a specific
Metal execution layout and its numerical behavior; CUDA source implementations
are related mechanisms, not measured baselines on the available hardware.

Reproducible summation has a substantial numerical-analysis literature.
Ahrens et al. construct order-independent accumulators \cite{repro}.
The implementation here instead retains the reference reduction behavior for
tested shapes on a pinned backend. TorchDynamo and TorchInductor reduce Python
and kernel overhead through graph compilation \cite{pytorch2}; the extension
shows whether gradient handling provides an additional benefit with that
execution path. The experiment compares these components directly rather than
attributing the compiler's gains to the gradient layout.

\section{Discussion and limitations}
@@DISCUSSION_TEXT@@
The saving is specific to the measured workload. Larger derivative matrix
products can reduce the fraction of time spent handling gradients, and more
distinct parameter lengths can reduce the opportunity for norm batching.
Conversely, dispatch overhead can be prominent when many small parameter
tensors are present. The study does not establish a hardware-independent
relationship between parameter count and speedup.

Numerical equality is also scoped. The real-arithmetic coordinate argument is
general under its assumptions, but the reduction matching depends on an actual
kernel implementation. Dense contiguous gradients, one hyperparameter group
and complete updates are enforced in this experiment. Shared or overlapping
parameter objects, sparse gradients, distributed synchronization, alternative
precision regimes and different backend versions need independent treatment.
Native embedding accumulation remains execution-sensitive; controlled equality
does not imply independently executed primary runs have identical final weights.

The experiment uses one Apple device and one reused corpus. V6 has twelve
pairs per configuration and descriptive quality intervals. Compiler artifacts
are warm-cached; cold compilation costs and other compiler/optimizer settings
remain uncharacterized. Both studies use per-step synchronization, which differs
from an asynchronously queued production training loop. The reported effect
should therefore be remeasured in the target deployment before extrapolation.
No conclusion is drawn about CUDA, distributed training, billion-parameter
models or model capability gains.

\section{Conclusion}
Cotangent batches gradient handling while retaining the native clipping
structure and complete model derivatives. On the measured eager MPS setup,
the independent replication meets its fixed time and quality requirements.
The extension identifies the contribution of individual components and tests
the method against executable compiled baselines on additional model layouts.
Its practical lesson is to check the actual reduction arithmetic when changing
tensor dispatch: mathematical equality alone does not certify numerical
equivalence, and optimizer fusion alone does not remove all post-backward cost.

\section*{Reproducibility}
The public repository contains the frozen protocols, complete raw curves,
per-step timings, source fingerprints, model/schedule hashes, analysis and this
standalone LaTeX source. Tags \code{protocol-v5} and \code{protocol-v6}
precede their official outcomes. Public records-only analysis recomputes
statistics without downloading data or weights; it explicitly leaves
checkpoint-file verification unperformed. All five earlier studies remain in
the experimental record, including their failed decisions. No observations are
removed to improve a reported result. @@RELEASE_LINK@@

\begin{thebibliography}{10}
\bibitem{adamw} I. Loshchilov and F. Hutter. Decoupled Weight Decay
Regularization. ICLR, 2019. \url{https://arxiv.org/abs/1711.05101}.
\bibitem{repro} W. Ahrens, J. Demmel, and H. D. Nguyen. Algorithms for
Efficient Reproducible Floating Point Summation. ACM Transactions on
Mathematical Software, 46(3), Article 22, 2020.
\url{https://doi.org/10.1145/3389360}.
\bibitem{pytorch2} J. Ansel et al. PyTorch 2: Faster Machine Learning Through
Dynamic Python Bytecode Transformation and Graph Compilation. ASPLOS, 2024.
\url{https://docs.pytorch.org/assets/pytorch2-2.pdf}.
\bibitem{fsdp} PyTorch contributors. FSDP flat parameter implementation.
\url{https://github.com/pytorch/pytorch/blob/main/torch/distributed/fsdp/_flat_param.py}.
\bibitem{deepspeed} DeepSpeed contributors. FusedAdam implementation.
\url{https://deepspeed.readthedocs.io/en/latest/_modules/deepspeed/ops/adam/fused_adam.html}.
\bibitem{apex} NVIDIA Apex contributors. Multi-tensor L2 norm kernel.
\url{https://github.com/NVIDIA/apex/blob/master/csrc/multi_tensor_l2norm_kernel.cu}.
\bibitem{mps} PyTorch contributors. MPS norm dispatcher, revision
\code{5c4886908584029761b579af026dcfb627c84070}.
\url{https://github.com/pytorch/pytorch/blob/5c4886908584029761b579af026dcfb627c84070/aten/src/ATen/native/mps/operations/ReduceOps.mm}.
\bibitem{wikitext} S. Merity, C. Xiong, J. Bradbury, and R. Socher. Pointer
Sentinel Mixture Models. ICLR, 2017. \url{https://arxiv.org/abs/1609.07843}.
\bibitem{gqa} J. Ainslie et al. GQA: Training Generalized Multi-Query Transformer
Models from Multi-Head Checkpoints. EMNLP, 2023. \url{https://arxiv.org/abs/2305.13245}.
\bibitem{glu} N. Shazeer. GLU Variants Improve Transformer. 2020.
\url{https://arxiv.org/abs/2002.05202}.
\bibitem{rms} B. Zhang and R. Sennrich. Root Mean Square Layer Normalization.
NeurIPS, 2019. \url{https://arxiv.org/abs/1910.07467}.
\bibitem{wavenet} A. van den Oord et al. WaveNet: A Generative Model for Raw
Audio. 2016. \url{https://arxiv.org/abs/1609.03499}.
\end{thebibliography}

\appendix
\section{Implementation outline}
\begin{enumerate}
\item Construct the model and attach its independent parameter leaves to
disjoint master-vector views before any forward pass.
\item Run the selected native or compiled forward/backward and materialize
every completed parameter gradient.
\item Group gradients by element count. Compute matched two-axis bucket norms,
using native scalar fallbacks for short or odd lengths.
\item Restore scalar norms to parameter order, compute the original joint norm,
and form the clipping coefficient.
\item Concatenate complete gradients, scale the vector, and apply native fused
AdamW with the original hyperparameters and scalar step count.
\end{enumerate}
The reference code is \code{research/v4/bucketed.py}. The extension's factorial
treatments and compiler composition are in \code{research/v6/methods.py}.
All prior source freezes remain unchanged.

\section{Experimental record}
\begin{table}[h]
\centering
\small
\begin{tabular}{llrl}
\toprule
Study & Method or purpose & Runs & Decision / scope\\
\midrule
V1 & Sampled weight gradients & 36 & Failed time and quality\\
V2 & Packing with flat norm & 24 & Failed combined gate\\
V3 & Deterministic adjoints, flat norm & 41 & Failed combined gate\\
V4 & Matched bucket norms & 70 & Failed quality uncertainty\\
V5 & Fixed-size independent replication & 200 & Passed fixed gates\\
V6 & Factorial and compiled extension & 128 & Descriptive comparisons\\
\bottomrule
\end{tabular}
\caption{Official planned training observations. Diagnostic pilots and numerical
verification steps are separate. V4 remains failed despite its timing advantage;
V5 is an independent fixed-size study of the unchanged method.}
\end{table}
The V4 upper 99\% quality bound was 0.01160035 BPB, above the original 0.01
margin. Its failure motivated power planning, not a relaxed threshold. Stronger
interval levels in later studies do not provide universal family-wise coverage
for every exploratory choice made during development.
\end{document}
